% संपूर्ण मराठी ग्रंथातील प्रकरण 44: मोडल चौकटींची परिभाष्यता
% हा संपादनयोग्य स्रोतखंड आहे. संकलनासाठी संपूर्ण स्रोतसंग्रहातील
% अक्षररूपे, पूर्वभाग, चिन्हव्याख्या आणि आकृती-स्रोत आवश्यक आहेत.
% मूळ: Open Logic Project; मजकूर आणि रूपांतर CC BY 4.0.
% यंत्रानुवाद, दुरुस्ती आणि तपासणी: OpenAI Codex —
% GPT-5.6 Sol आणि GPT-6 Sol, दोन्ही Ultra effort.
% दुरुस्ती-पुनरावलोकन: GPT-6.1 Sol, Ultra effort.
\chapter{मोडल चौकटींची परिभाष्यता}\label{nml:frd:chap}









\section{प्रस्तावना}\label{nml:frd:int:sec}

मोडल तर्कशास्त्रात प्राप्यता संबंध आणि तो संबंध असलेल्या
प्रतिमानांतील काही सूत्रांची सत्यता यांमधील संबंध
महत्त्वाचा असतो. उदाहरणार्थ, प्राप्यता संबंध परावर्ती
आहे, म्हणजे प्रत्येक $w
\in W$ साठी $Rww$ आहे, असे समजा. दुसऱ्या शब्दांत,
प्रत्येक जग स्वतःपासून प्राप्य आहे. मग एखाद्या जगात~$w$
येथे $\Box \varphi$ सत्य असेल, तर ज्या प्राप्य जगांत~$\varphi$
सत्य असणे आवश्यक आहे त्यांत $w$ स्वतःही येते. म्हणून
$\mModel{M}$ चा प्राप्यता संबंध~$R$ परावर्ती असेल, तर
कोणतेही जग $w$ आणि कोणतेही सूत्र $\varphi$ घेतले, तरी
$\Box \varphi
\lif \varphi$ तेथे सत्य असेल. दुसऱ्या शब्दांत, $\Box p \lif
p$ हा रूपबंध आणि त्याचे सर्व आदेशन-दृष्टांत~$\mModel{M}$
मध्ये सत्य आहेत.

मात्र व्यत्यास असत्य आहे. उदाहरणार्थ, $\Box p \lif p$
हे~$\mModel{M}$ मध्ये सत्य असल्यावरून $R$ परावर्ती आहे
असे ठरत नाही. कारण $\Box p \lif
p$ सर्व जगांत सत्य असलेले अपरावर्ती प्रतिमान~$\mModel{M}$
सहज मिळते: स्वतःपासून प्राप्य नसलेले एकच जग~$w$ असलेले,
पण $w \in V(p)$ असलेले प्रतिमान घ्या. $p$ चे सत्यतामूल्य
योग्य निवडल्यास अपरावर्ती प्रतिमानात एखाद्या विशिष्ट
सूत्रासाठी $\Box \varphi \lif \varphi$ सत्य करता येते.

यासाठी मूल्यांकनाचा घटक~$V$ बाजूला ठेवू. $V(p)$ मध्ये
कोणती जगं आहेत याची पर्वा न करता, $\Box p \lif p$ हे
$\mModel{M}$ मधील सर्व जगांत सत्य असावे अशी अट घातल्यास
$R$ परावर्ती असणे आवश्यक ठरते. कारण अपरावर्ती
प्रतिमानात $Rww$ नसलेले किमान एक जग~$w$ असते.
$V(p) =
W \setminus \{w\}$ असे ठेवा. मग~$w$ सोडून इतर
सर्व जगांत $p$ सत्य असते; म्हणून~$w$ पासून प्राप्य
असलेल्या सर्व जगांतही ते सत्य असते ($w$ हे~$w$ पासून
प्राप्य नसल्याची खात्री आहे, आणि $p$ असत्य असलेले
$w$ हेच एकमेव जग आहे). दुसरीकडे, $w$ येथे $p$ असत्य
आहे; म्हणून $\Box
p \lif p$ हे~$w$ येथे असत्य आहे.

यावरून मूल्यांकनाशिवाय प्रतिमानाची रचना दाखवण्यासाठी
स्वतंत्र संकेतन हवे, असे सुचते; अशा रचनांना \emph{चौकटी}
म्हणतो. चौकट $\mModel{F}$ म्हणजे जगांचा संच आणि
प्राप्यता संबंध यांची जोडी $\tuple{W, R}$.
प्रत्येक प्रतिमान $\tuple{W, R, V}$ हे $\tuple{W, R}$
या चौकटीवर \emph{आधारित} असते. उलट, एखादी चौकट
तिच्यावर आधारित प्रतिमानांचा वर्ग ठरवते; आणि चौकटींचा
वर्ग त्यांपैकी कोणत्यातरी चौकटीवर आधारित प्रतिमानांचा
वर्ग ठरवतो. सूत्राचे चौकटीत \emph{वैध} असणे,
म्हणजे $\mModel{F} \Entails \varphi$, अशी व्याख्या करता येते:
चौकट~$\mModel{F}$ वर आधारित प्रत्येक $\mModel{M}$
साठी $\mSat{M}{\varphi}$.

या संकेतनाने सूत्रे आणि चौकटींचे वर्ग यांतील
अनुरूपता सांगता येते: उदाहरणार्थ,
$\mModel{F} \Entails \Box p \lif p$ तेव्हा आणि केवळ
तेव्हाच, जेव्हा $\mModel{F}$ परावर्ती आहे.













\section{प्राप्यता संबंधांचे गुणधर्म}\label{nml:frd:acc:sec}

अनेक मोडल सूत्रे प्राप्यता संबंधाच्या साध्या, कधीकधी
परिचित, गुणधर्मांचे वैशिष्ट्य दर्शवतात. एका दिशेने याचा
अर्थ असा की, एखादा गुणधर्म असलेल्या प्रत्येक प्रतिमानात
संबंधित सूत्र आणि त्याचे सर्व आदेशन-दृष्टांत सत्य
असतात. प्राप्यता संबंधांचे पाच अभिजात प्रकार आणि त्यांनी
सत्य ठरवलेली सूत्रे प्रथम पाहू.

\begin{thm}\label{nml:frd:acc:thm:soundschemas}
  $\mModel{M} = \tuple{W, R, V}$ हे प्रतिमान असू द्या.
  $R$ ला \ref{nml:frd:acc:tab:five} च्या डाव्या बाजूचा गुणधर्म
  असेल, तर उजव्या बाजूच्या सूत्राचा प्रत्येक
  दृष्टांत~$\mModel{M}$ मध्ये सत्य असतो.
\end{thm}

\begin{table}[t]
    \begin{tabular}{| l || l |}
      \hline
      {\emph{$R$ पुढील प्रकारचा असेल \dots}} & {\emph{तर \dots सूत्र~$\mModel{M}$ मध्ये सत्य:}} \\
      \hline \hline
      \emph{सर्वत्र उत्तरवर्ती}: $\forall u \exists v Ruv$ & \hbox to.43\textwidth
           {$\Box p \lif \Diamond p$ \hfill (\Ax{D})} \\
      \hline
      \emph{परावर्ती}: $\forall w Rww$
      & \hbox to.43\textwidth
          {$\Box p \lif p$ \hfill (\Ax{T})} \\
      \hline
      \emph{सममित}: &  \hbox to .43\textwidth
           {$p \lif \Box\Diamond p$ \hfill (\Ax{B})} \\
           $\forall u\forall v(Ruv \lif Rvu)$ & \\
      \hline
      \emph{संक्रमक}: & \hbox to.43\textwidth
           {$\Box p \lif \Box \Box p$ \hfill (\Ax{4})} \\
      $\forall u \forall v \forall w ((Ruv \land Rvw) \lif Ruw)$ & \\
      \hline
      \emph{युक्लिडी}: & \hbox to .43\textwidth
        {$\Diamond p \lif \Box\Diamond p$ \hfill (\Ax{5})} \\
      $\forall w \forall u \forall v ((Rwu \land Rwv) \lif Ruv)$ &\\
      \hline
    \end{tabular}
    \caption{अनुरूपतेविषयी पाच तथ्ये.}
    \label{nml:frd:acc:tab:five}
  \end{table}


\begin{proof}
  \Ax{B} चे प्रकरण पाहू. रूपबंध प्रतिमानात सत्य आहे हे
  दाखवण्यासाठी त्याचे सर्व दृष्टांत त्या प्रतिमानातील
  प्रत्येक जगात सत्य आहेत हे दाखवावे लागते. म्हणून
  $\varphi \lif \Box\Diamond \varphi$ हा \Ax{B} चा कोणताही
  दृष्टांत, आणि $w \in W$ हे कोणतेही जग असू द्या.
  $\Box \Diamond
  \varphi$ हे $w$ येथे सत्य असल्याचे दाखवण्यासाठी पूर्वांग
  $\varphi$ हे $w$ येथे सत्य आहे असे समजा. मग $w$ पासून
  प्राप्य असलेल्या प्रत्येक $w'$ येथे $\Diamond \varphi$
  सत्य आहे हे दाखवायचे आहे. आता $Rww'$ असलेल्या
  कोणत्याही $w'$ साठी, सममितीच्या गृहीतकावरून $Rw'w$
  देखील असते (\ref{nml:frd:acc:fig:Bsymm} पाहा).
  $\mSat{M}{\varphi}[w]$ असल्यामुळे
  $\mSat{M}{\Diamond \varphi}[w']$. $Rww'$ असलेले $w'$
  कोणतेही असल्यामुळे $\mSat{M}{\Box\Diamond\varphi}[w]$.

  इतर प्रकरणे सरावासाठी ठेवली आहेत.
\end{proof}

\begin{prob}
  \ref{nml:frd:acc:thm:soundschemas} ची सिद्धता पूर्ण करा.
\end{prob}

\begin{figure}
  \begin{center}
    \begin{tikzpicture}[modal]
      \node[world] (w1) [label={below:$\mSat{{}}{\formula{A}}$\\
          $\mSat{{}}{\Box\Diamond \formula{A}}$}] {$w$};
      \node[world] (w2) [label=below:{$\mSat{{}}{\Diamond \formula{A}}$},
        right=of w1] {$w'$};
      \draw[->, bend left] (w1) to (w2);
      \draw[->, bend left] (w2) to (w1);
    \end{tikzpicture}
  \end{center}
\caption{सममितीवर आधारित युक्तिवाद.}
\label{nml:frd:acc:fig:Bsymm}
\end{figure}

\ref{nml:frd:acc:thm:soundschemas} मधील व्यत्यास खरे नाहीत, हे लक्षात
घ्या: एखादा रूपबंध प्रतिमानात सत्य आहे म्हणून त्या प्रतिमानाच्या
प्राप्यता संबंधाला संबंधित गुणधर्म असतोच असे नाही.
\Ax{T} आणि परावर्ती प्रतिमाने या प्रकरणात, \Ax{T} स्वतःच
असत्य असलेले प्रतिमान सहज देता येते: $W = \{w\}$ आणि
$V(p) = \emptyset$ घ्या; प्राप्यता संबंध रिकामा असू द्या.
मग $R$ परावर्ती नाही, पण $\mSat{M}{\Box
  p}[w]$ आणि $\mSat/{M}{p}[w]$. तरी येथे~$\mModel{M}$
मध्ये \Ax{T} चा असत्य असलेला एक दृष्टांत आपण दाखवला आहे; काही इतर दृष्टांत,
उदाहरणार्थ $\Box \lnot p \lif \lnot p$, सत्य आहेत.
\Ax{T} चे \emph{प्रत्येक आदेशन-दृष्टांत}~$\mModel{M}$
मध्ये सत्य असूनही $\mModel{M}$ परावर्ती नसते, असे
प्रतिमान शोधणे अधिक कठीण आहे. पण अशी प्रतिमानेही आहेत:

\begin{prop}\label{nml:frd:acc:prop:reflexive}
  $\mModel{M} = \tuple{W, R, V}$ हे प्रतिमान असू द्या,
  जिथे $W = \{u, v
  \}$, जगं $u$ आणि $v$ यांचा $R$ संबंध दोन्ही दिशांनी
  आहे, म्हणजे $Ruv$ आणि $Rvu$, आणि कोणतेही जग स्वतःशी
  संबंधित नाही. प्रत्येक $p$ साठी $u \in V(p) \Leftrightarrow v
  \in V(p)$ असे समजा. मग:
  \begin{enumerate}
  \item प्रत्येक $\varphi$ साठी: $\mSat{M}{\varphi}[u]$ तेव्हा आणि
    केवळ तेव्हाच, जेव्हा $\mSat{M}{\varphi}[v]$
    ($\varphi$ वर विगमन वापरा).
  \item \Ax{T} चा प्रत्येक दृष्टांत $\mModel{M}$ मध्ये सत्य आहे.
  \end{enumerate}
  $\mModel{M}$ परावर्ती नाही (खरेतर ते
  \emph{अपरावर्ती} आहे), म्हणून \Ax{T} च्या बाबतीत
  \ref{nml:frd:acc:thm:soundschemas} चा व्यत्यास असत्य आहे.
  \ref{nml:frd:acc:thm:soundschemas} मधील काही, पण सर्वच नव्हे,
  इतर रूपबंधांसाठीही असेच युक्तिवाद देता येतात.
\end{prop}

\begin{prob}
  \ref{nml:frd:acc:prop:reflexive} मधील दावे सिद्ध करा.
\end{prob}

आपण \Ax{D}, \Ax{T}, \Ax{B}, \Ax{4}, आणि \Ax{5}
या पाच अभिजात सूत्रे वर भर देणार असलो, तरी
\ref{nml:frd:acc:tab:anotherfive} मध्ये प्राप्यता संबंधांचे आणखी
काही गुणधर्म नोंदवले आहेत. प्रत्येक जगापासून जास्तीत
जास्त एक जग प्राप्य असेल, तर प्राप्यता संबंध~$R$
आंशिक फलनीय असतो. प्रत्येक जगापासून नेमके
एकच जग प्राप्य असेल, तर त्याला फलनीय म्हणतो.
(म्हणून फलनीय संबंध नेमके तेच आहेत जे सर्वत्र उत्तरवर्ती
आणि आंशिक फलनीय दोन्ही आहेत.) त्यांना ``फलनीय''
म्हणतात कारण प्राप्यता संबंध एखाद्या (आंशिक) फलनासारखा
वागतो. $Ruv$ असेल तेव्हा $u$ आणि~$v$ यांच्या
``मध्ये'' एखादे~$w$ असेल, तर संबंध दुर्बल सघन
असतो. एका अर्थाने दुर्बल सघन संबंध संक्रमक संबंधांच्या
उलट आहेत: संक्रमक संबंधात~$u$ पासून~$w$ मार्गे
$v$ पर्यंत जाता येत असेल, तर $u$ पासून $v$ पर्यंत
थेट जाता येते; दुर्बल सघन संबंधात~$u$ पासून $v$
पर्यंत थेट जाता येत असेल, तर एखाद्या~$w$ मार्गेही
जाता येते. एखाद्या जग~$w$ पासून जगं $u$ आणि~$v$
प्राप्य असतील तेव्हा $u$ आणि~$v$ या दोन्हींपासून
एकच जग~$t$ प्राप्य असेल, तर संबंध दुर्बल दिशित असतो;
यालाच कधीकधी ``हिऱ्याचा गुणधर्म'' किंवा ``संगमन''
म्हणतात.

\begin{table}[t]
    \setlength{\tabcolsep}{4pt}
\begin{tabular}{| p{.50\textwidth} || p{.44\textwidth} |}
      \hline
      {\emph{$R$ पुढील प्रकारचा असेल \dots}} & {\emph{तर \dots सूत्र~$\mModel{M}$ मध्ये सत्य:}} \\
      \hline \hline
      \emph{आंशिक फलनीय}: &
      \multirow{2}{*}{$\Diamond p \lif \Box p$} \\
      $ \forall w \forall u \forall v ((Rwu \land Rwv) \lif u =v )$ & \\
      \hline
      \multirow{2}{*}{\emph{फलनीय}: $\forall w \exists v \forall u(Rwu \liff \eq[u][v]) $}
      & \multirow{2}{*}{$\Diamond p \liff \Box p$} \\
      & \\
      \hline
      \emph{दुर्बल सघन}: &  \multirow{2}{*}{$\Box \Box p \lif
        \Box p$} \\
      $\forall u\forall v(Ruv \lif \exists w(Ruw \land Rwv))$ & \\
      \hline
      \emph{दुर्बल संयोजित}: &
      \multirow{3}{*}{\hbox to.43\textwidth{$
        \begin{array}{@{}l@{}}
          \Box ((p \land \Box p) \lif q) \lor {} \\
          \qquad \Box ((q \land \Box q) \lif p)
        \end{array}$ \hfill (\Ax{L})}
      } \\
      $\forall w \forall u \forall v ((Rwu \land Rwv) \lif {}$ &\\
      \qquad $(Ruv \lor u=v \lor Rvu))$ & \\
      \hline
      \emph{दुर्बल दिशित}: & \multirow{3}{*}{\hbox to .43\textwidth
        {$\Diamond\Box p \lif \Box\Diamond p$\hfill (\Ax{G})}} \\
      $\forall w \forall u \forall v ((Rwu \land Rwv) \lif {}$ & \\
      \qquad $\exists t (Rut \land Rvt))$ & \\
      \hline
    \end{tabular}
    \caption{अनुरूपतेविषयी आणखी पाच तथ्ये.}
    \label{nml:frd:acc:tab:anotherfive}
  \end{table}

\begin{prob}
  $\mModel{M} = \tuple{W, R, V}$ हे प्रतिमान असू द्या.
  $R$ ला \ref{nml:frd:acc:tab:anotherfive} मधील
  डावीकडील गुणधर्म असेल, तर उजवीकडील संबंधित सूत्राचा
  प्रत्येक दृष्टांत~$\mModel{M}$ मध्ये सत्य असतो, हे दाखवा.
\end{prob}














\section{चौकटी}\label{nml:frd:fra:sec}

\begin{defn}
  \emph{चौकट} म्हणजे $\mModel{F} = \tuple{W,R}$ ही जोडी, जिथे
  $W$ हा जगांचा रिक्त नसलेला संच आणि $R$ हा~$W$ वरील
  द्विस्थानी संबंध आहे. प्रतिमान~$\mModel{M}$ हे
  $\mModel{F} = \tuple{W,R}$ या चौकटीवर \emph{आधारित}
  आहे तेव्हा आणि केवळ तेव्हाच, जेव्हा काही मूल्यांकन
  $V$ साठी $\mModel{M} = \tuple{W, R, V}$.
\end{defn}

\begin{defn}
  $\mModel{F}$ ही चौकट असेल, तर~$\mModel{F}$ वर आधारित
  प्रत्येक प्रतिमान~$\mModel{M}$ साठी $\mSat{M}{\varphi}$
  असल्यास $\varphi$ हे \emph{$\mModel{F}$ मध्ये वैध,}
  म्हणजे $\mModel{F} \Entails \varphi$, असे म्हणतो.

  $\mClass{F}$ हा चौकटींचा वर्ग असेल, तर त्या वर्गातील
  प्रत्येक चौकट~$\mModel{F} \in \mClass{F}$ साठी
  $\mModel{F} \Entails \varphi$ असल्यास, आणि तेव्हाच,
  $\varphi$ हे \emph{$\mClass{F}$ मध्ये वैध,}
  म्हणजे $\mClass{F} \Entails \varphi$, असे म्हणतो.
\end{defn}

चौकटी महत्त्वाच्या आहेत, कारण रूपबंध आणि प्राप्यता
संबंध~$R$ याचे गुणधर्म यांमधील अनुरूपता \emph{प्रतिमानांच्या
नव्हे, तर चौकटींच्या} पातळीवर असते. उदाहरणार्थ,
\Ax{T} सर्व परावर्ती प्रतिमानांत सत्य असला,
तरी \Ax{T} सत्य असलेल्या प्रत्येक प्रतिमानाचा संबंध
परावर्ती असतोच असे नाही. परंतु असे \emph{खरे} आहे की,
\Ax{T} सर्व परावर्ती \emph{चौकटींवर} \emph{वैध} आहे,
आणि \Ax{T} वैध
असलेली प्रत्येक चौकट परावर्ती आहे, ही दोन्ही विधाने
खरी आहेत.

\begin{rem}
चौकटींच्या वर्गातील वैधता ही प्रतिमानांच्या वर्गातील
वैधतेच्या संकल्पनेचे विशेष प्रकरण आहे: $\mClass{F} \Entails \varphi$
तेव्हा आणि केवळ तेव्हाच, जेव्हा $\mClass{C} \Entails \varphi$,
जिथे $\mClass{C}$ हा~$\mClass{F}$ मधील एखाद्या
चौकटीवर आधारित सर्व प्रतिमानांचा वर्ग आहे.

स्पष्टच आहे की, सूत्र किंवा रूपबंध वैध असेल,
म्हणजे \emph{सर्व} प्रतिमानांच्या वर्गाच्या संदर्भात वैध
असेल, तर तो चौकटींच्या कोणत्याही वर्ग~$\mClass{F}$
च्या संदर्भातही वैध असतो.
\end{rem}













\section{चौकटींची परिभाष्यता}\label{nml:frd:def:sec}

\ref{nml:frd:acc:thm:soundschemas} मधील व्यत्यास प्रतिमानांसाठी
असत्य असले, तरी ``प्रतिमान'' याऐवजी ``चौकट'' घेतल्यास
ते खरे ठरतात: \ref{nml:frd:acc:thm:soundschemas} मध्ये पाहिलेल्या
गुणधर्मांसाठी, सूत्र एखाद्या \emph{चौकटीत} वैध
असेल, तर त्या चौकटीच्या प्राप्यता संबंधाला संबंधित गुणधर्म
\emph{असतोच}. म्हणून संबंधित गुणधर्म असलेल्या चौकटींचे
वर्ग हे सूत्रे \emph{परिभाषित करतात}.

\begin{defn}
  $\mClass{F}$ हा चौकटींचा वर्ग असेल, तर नेमक्या
  त्यातील चौकटी~$\mModel{F} \in \mClass{F}$
  आणि त्यांच्याच बाबतीत $\mModel{F} \Entails \varphi$ असेल,
  तर $\varphi$ हा $\mClass{F}$ वर्ग \emph{परिभाषित करतो},
  असे म्हणतो.
\end{defn}

आता चौकटींसाठी परिभाष्यतेचे पूर्ण निष्कर्ष सिद्ध करू.

\begin{thm}\label{nml:frd:def:thm:fullCorrespondence}
\ref{nml:frd:acc:tab:five} च्या उजव्या बाजूचे सूत्र
चौकट~$\mModel{F}$ मध्ये वैध असेल, तर $\mModel{F}$ ला
डाव्या बाजूचा गुणधर्म असतो.
\end{thm}

\begin{proof}
  \begin{enumerate}
  \item \Ax{D} हे $\mModel{F} = \tuple{W, R}$ मध्ये वैध आहे,
    म्हणजे $\mModel{F} \Entails \Box p \lif \Diamond p$,
    असे समजा. $\mModel{F}$ वर आधारित
    $\mModel{M} = \tuple{W, R, V}$ हे प्रतिमान आणि
    $w \in W$ घ्या. $Rwv$ असलेले एखादे $v$ आहे हे
    दाखवायचे आहे. तसे नाही असे समजा. मग~$p$ सहित
    कोणत्याही~$\varphi$ साठी $\mSat{M}{\Box \varphi}[w]$ आणि
    $\mSat/{M}{\Diamond \varphi}[w]$ ही दोन्ही विधाने खरी
    असतात. पण मग $\mSat/{M}{\Box p \lif
      \Diamond p}[w]$, आणि हे $\mModel{F}
    \Entails \Box p \lif \Diamond p$ या गृहीतकाला
    विरोधी आहे.
  \item \Ax{T} हे $\mModel{F}$ मध्ये वैध आहे,
    म्हणजे $\mModel{F} \Entails \Box
    p \lif p$, असे समजा. कोणतेही जग $w \in W$ घ्या;
    $Rww$ हे दाखवायचे आहे. $u \in V(p)$ तेव्हा आणि
    केवळ तेव्हाच, जेव्हा $Rwu$, असे मूल्यांकन
    ठरवा (आणि $p$ पेक्षा वेगळ्या $q$ साठी $V(q)$
    मनमाना ठेवा, उदाहरणार्थ $V(q) = \emptyset$).
    $\mModel{M} = \tuple{W, R, V}$ घ्या. या रचनेनुसार
    $Rwu$ असलेल्या सर्व $u$ साठी $\mSat{M}{p}[u]$;
    म्हणून $\mSat{M}{\Box p}[w]$. पण गृहीतकावरून
    $\Box p \lif p$ हे~$w$ येथे सत्य आहे; म्हणून
    $\mSat{M}{p}[w]$. $V$ च्या व्याख्येनुसार हे
    $Rww$ असेल तेव्हाच शक्य आहे.
  \item प्रतिपरिवर्तनाने सिद्ध करू: $\mModel{F}$ सममित
    नाही असे समजा. मग \Ax{B}, म्हणजे
    $p \lif \Box\Diamond p$, हे
    $\mModel{F}= \tuple{W, R}$ मध्ये वैध नाही हे
    दाखवू. $\mModel{F}$ सममित नसल्यामुळे
    $u$, $v \in W$ अशी जगं आहेत की $Ruv$ आहे,
    पण $Rvu$ नाही. $w \in V(p)$ तेव्हा आणि केवळ
    तेव्हाच, जेव्हा $Rvw$ नाही, असे $V$ ठरवा (इतरत्र
    $V$ मनमाना आहे). $\mModel{M} =\tuple{W, R,
      V}$ घ्या. $V$ च्या व्याख्येनुसार, $Rvw$
    नसलेल्या सर्व $w$ साठी $\mSat{M}{p}[w]$;
    विशेषतः $Rvu$ नसल्यामुळे $\mSat{M}{p}[u]$.
    तसेच $Rvw$ तेव्हा आणि केवळ तेव्हाच, जेव्हा
    $w \notin V(p)$; म्हणून $Rvw$ आणि
    $\mSat{M}{p}[w]$ या दोन्ही अटी पूर्ण करणारे
    कोणतेही~$w$ नाही. त्यामुळे
    $\mSat/{M}{\Diamond p}[v]$. $Ruv$ असल्याने
    $\mSat/{M}{\Box\Diamond p}[u]$ देखील. म्हणून
    $\mSat/{M}{p \lif
      \Box\Diamond p}[u]$, आणि त्यामुळे
    $\Ax{B}$ हे~$\mModel{F}$ मध्ये वैध नाही.
  \item \Ax{4} हे $\mModel{F} = \tuple{W,R}$ मध्ये
    वैध आहे, म्हणजे $\mModel{F} \Entails \Box p
    \lif \Box\Box p$, असे समजा. $u$, $v$, $w \in W$
    ही कोणतीही जगं $Ruv$ आणि $Rvw$ अशा अटी पूर्ण
    करत असू द्या; $Ruw$ सिद्ध करायचे आहे.
    $z \in V(p)$ तेव्हा आणि केवळ तेव्हाच, जेव्हा
    $Ruz$, असे $V$ ठरवा (इतरत्र $V$ मनमाना आहे).
    $\mModel{M} =\tuple{W, R, V}$ घ्या. $V$ च्या
    व्याख्येनुसार $Ruz$ असलेल्या सर्व $z$ साठी
    $\mSat{M}{p}[z]$; म्हणून $\mSat{M}{\Box p}[u]$.
    गृहीतकावरून \Ax{4}, म्हणजे $\Box p
    \lif \Box \Box p$, हे $u$ येथे सत्य आहे;
    म्हणून $\mSat{M}{\Box \Box
      p}[u]$. $Ruv$ आणि $Rvw$ असल्यामुळे
    $\mSat{M}{p}[w]$; पण $V$ च्या व्याख्येनुसार
    हे $Ruw$ असेल तेव्हाच शक्य आहे.
  \item पुन्हा प्रतिपरिवर्तन वापरू. चौकट
    $\mModel{F} = \tuple{W, R}$ युक्लिडी नाही असे
    मानून ती \Ax{5} असत्य ठरवते, म्हणजे
    $\mModel{F} \Entails/ \Diamond p \lif
      \Box\Diamond p$, हे दाखवू. $u$, $v$, $w \in W$
    अशी जगं असू द्या की $Rwu$ आणि $Rwv$ आहेत,
    पण $Ruv$ नाही. सर्व जगं $z$ साठी $z \in V(p)$
    तेव्हा आणि केवळ तेव्हाच, जेव्हा $Ruz$ असे
    \emph{नाही}, असे $V$ ठरवा. $\mModel{M}
    =\tuple{W, R, V}$ घ्या. मग गृहीतकावरून
    $\mSat{M}{p}[v]$, आणि $Rwv$ असल्यामुळे
    $\mSat{M}{\Diamond p}[w]$ देखील. परंतु
    $Ruy$ आणि $\mSat{M}{p}[y]$ असलेले कोणतेही
    जग~$y$ नाही; म्हणून $\mSat/{M}{\Diamond
      p}[u]$. $Rwu$ असल्यामुळे
    $\mSat/{M}{\Box\Diamond
      p}[w]$; म्हणून \Ax{5}, म्हणजे
    $\Diamond p \lif \Box\Diamond p$, हे~$w$
    येथे असत्य आहे.
  \end{enumerate}
\end{proof}

\Ax{D} ची सिद्धता आणि इतर प्रकरणे यांत एक फरक दिसेल:
त्यात मूल्यांकन~$V$ ची कोणतीही विशेष निवड करावी लागली नाही. प्रत्यक्षात,
$\mSat{M}{\Ax{D}}$ असल्यास $\mModel{M}$ सर्वत्र
उत्तरवर्ती आहे, हे आपण सिद्ध केले. म्हणून \Ax{D}
हा रूपबंध केवळ चौकटींचा नव्हे, तर सर्वत्र उत्तरवर्ती
\emph{प्रतिमानांचा} वर्ग परिभाषित करतो.

\begin{cor}\label{nml:frd:def:prop:D-serial}
  \Ax{D} सत्य असलेले कोणतेही प्रतिमान सर्वत्र उत्तरवर्ती असते.
\end{cor}

\begin{cor}
\ref{nml:frd:acc:tab:five} च्या उजव्या बाजूचे प्रत्येक
सूत्र डावीकडील गुणधर्म असलेल्या चौकटींचा
वर्ग परिभाषित करते.
\end{cor}

\begin{proof}
  \ref{nml:frd:acc:thm:soundschemas} मध्ये आपण सिद्ध केले की,
  प्रतिमानाला डावीकडील गुणधर्म असेल, तर उजवीकडचे
  सूत्र त्यात सत्य असते. म्हणून चौकट~$\mModel{F}$
  ला तो गुणधर्म असेल, तर उजवीकडचे सूत्र
  $\mModel{F}$ मध्ये वैध असते.
  \ref{nml:frd:def:thm:fullCorrespondence} मध्ये आपण व्यत्यास
  सिद्ध केले: उजवीकडचे सूत्र~$\mModel{F}$
  मध्ये वैध असेल, तर $\mModel{F}$ ला डावीकडचा
  गुणधर्म असतो.
\end{proof}

\begin{prob}
\ref{nml:frd:acc:tab:anotherfive} च्या उजव्या
बाजूचे सूत्र चौकट~$\mModel{F}$ मध्ये वैध असेल,
तर $\mModel{F}$ ला डावीकडील गुणधर्म असतो, हे दाखवा.
यासाठी डावीकडील गुणधर्म \emph{नसलेली} चौकट घ्या,
आणि उजवीकडचे सूत्र एखाद्या जगात असत्य ठरेल
असे योग्य~$V$ ठरवा.
\end{prob}

\ref{nml:frd:def:thm:fullCorrespondence} वरून गुणधर्म एकत्रही
करता येतात हे दिसते: उदाहरणार्थ \Ax{B} आणि
\Ax{4} हे दोन्ही~$\mModel{F}$ मध्ये वैध असतील,
तर चौकट सममित आणि संक्रमक दोन्ही असते.
अनेक महत्त्वाची मोडल तर्कशास्त्रे काही चौकट-गुणधर्म
एकत्र असलेल्या सर्व चौकटींमध्ये वैध असलेल्या
सूत्रांचे संच म्हणून निश्चित होतात. म्हणून
त्यांना संबंधित परिभाषक सूत्रे वैध असलेल्या
सर्व चौकटींमध्ये वैध असलेल्या सूत्रांचे संच
म्हणूनही सांगता येते. उदाहरणार्थ, अभिजात व्यवस्था
\Log{S4} म्हणजे सर्व परावर्ती आणि संक्रमक चौकटींमध्ये
वैध असलेल्या सूत्रांचा संच; म्हणजे \Ax{T}
आणि~\Ax{4} दोन्ही वैध असलेल्या सर्व चौकटींमध्ये
वैध असलेल्या सूत्रांचा संच.
\Log{S5} म्हणजे सर्व परावर्ती, सममित आणि युक्लिडी
चौकटींमध्ये वैध असलेल्या सूत्रांचा संच; म्हणजे
\Ax{T}, \Ax{B}, आणि \Ax{5} ही तिन्ही वैध
असलेल्या सर्व चौकटींमध्ये वैध असलेल्या सूत्रांचा संच.

$R$ च्या गुणधर्मांमधील तार्किक संबंध सामान्यतः
संबंधित परिभाषक सूत्रे मधील संबंधांशी जुळतात.
उदाहरणार्थ, प्रत्येक परावर्ती संबंध सर्वत्र उत्तरवर्ती
असतो; म्हणून चौकटीत \Ax{T} वैध असेल, तर
\Ax{D} देखील वैध असतो. (हा संबंध \emph{तार्किक
निष्पन्नतेचा नाही}, हे लक्षात घ्या. $\mSat{M}{\Ax{T}}[w]$
असल्यावर नेहमी $\mSat{M}{\Ax{D}}[w]$ असते, असे
नाही.) असे काही संबंध पुढे नोंदवतो.

\begin{prop}\label{nml:frd:def:prop:relation-facts}
  $R$ हा $W$ या संचावरील द्विस्थानी संबंध असू द्या. मग:
  \begin{enumerate}
  \item $R$ परावर्ती असेल, तर तो सर्वत्र उत्तरवर्ती असतो.
  \item $R$ सममित असेल, तर तो संक्रमक असतो तेव्हा आणि
    केवळ तेव्हाच, जेव्हा तो युक्लिडी असतो.
  \item $R$ सममित किंवा युक्लिडी असेल, तर तो दुर्बल
    दिशित असतो (त्याला ``हिऱ्याचा गुणधर्म'' असतो).
  \item $R$ युक्लिडी असेल, तर तो दुर्बल संयोजित असतो.
  \item $R$ फलनीय असेल, तर तो सर्वत्र उत्तरवर्ती असतो.
  \end{enumerate}
\end{prop}

\begin{prob}
  \ref{nml:frd:def:prop:relation-facts} सिद्ध करा.
\end{prob}












\section{प्रथम-क्रमातील परिभाष्यता}\label{nml:frd:fol:sec}

चौकटींच्या प्राप्यता संबंधांचे अनेक गुणधर्म मोडल
सूत्रे नी परिभाषित करता येतात, हे आपण पाहिले.
उदाहरणार्थ, चौकटींची सममिती \Ax{B}, म्हणजे
$p \lif \Box\Diamond
p$, या सूत्र ने परिभाषित होते. आतापर्यंत
पाहिलेल्या सर्व अटी एकच द्विस्थानी विधेय
असलेल्या भाषा तील प्रथम-क्रम सूत्रे
नी व्यक्त करता येतात. उदाहरणार्थ, सममिती
$\lforall[x][\lforall[y][(\Atom{Q}{x,y} \lif \Atom{Q}{y, x})]]$
याने परिभाषित होते: $\Domain{M} =
W$ आणि $\Assign{Q}{M} = R$ असलेली प्रथम-क्रम
संरचना~$\Struct{M}$ वरील सूत्राची पूर्ति करते
तेव्हा आणि केवळ तेव्हाच, जेव्हा $R$ सममित असतो.
यातून पुढील व्याख्या सुचते:

\begin{defn}
  चौकटींचा वर्ग~$\mClass{F}$ \emph{प्रथम-क्रमात
  परिभाष्य} असतो, जर एकच द्विस्थानी विधेय~$Q$
  असलेल्या प्रथम-क्रम भाषेतील असे वाक्य~$\varphi$
  असेल की, $\Domain{M} = W$ आणि $\Assign{Q}{M}
  = R$ असलेल्या प्रथम-क्रम संरचना~$\Struct{M}$
  मध्ये $\Sat{M}{\varphi}$ तेव्हा आणि केवळ तेव्हाच,
  जेव्हा $\mModel{F} =
  \tuple{W, R} \in \mClass{F}$.
\end{defn}

आतापर्यंत पाहिलेले गुणधर्म आणि त्यांना परिभाषित करणारी
मोडल सूत्रे ही अपवादात्मक उदाहरणे ठरतात.
प्रत्येक सूत्र प्रथम-क्रमात परिभाष्य चौकटींचाच
वर्ग परिभाषित करते असे नाही; तसेच प्रथम-क्रमात
परिभाष्य चौकटींचा प्रत्येक वर्ग मोडल सूत्राने
परिभाषित होतो असेही नाही.

पहिल्या दाव्याचे प्रतिउदाहरण L\"ob सूत्र देते:
\begin{equation}
\Box(\Box p \lif p) \lif \Box p. \tag{\Ax{W}}
\end{equation}
\Ax{W} संक्रमक आणि प्रतिलोम-सुस्थापित चौकटींचा वर्ग
परिभाषित करते. $Rw_2w_1$, $Rw_3w_2$, \dots{}
अशी $w_1$, $w_2$, \dots{} ही अनंत श्रेणी नसल्यास
संबंध सुस्थापित असतो. उदाहरणार्थ, $\Nat$ वरील
$<$ हा संबंध सुस्थापित आहे; पण $\Int$ वरील
$<$ सुस्थापित नाही. एखाद्या संबंधाचा प्रतिलोम
सुस्थापित असेल, तर आणि तेव्हाच, तो संबंध प्रतिलोम-सुस्थापित
असतो. म्हणजे प्रतिलोम-सुस्थापित संबंधात $Rw_1w_2$,
$Rw_2w_3$, \dots{} अशी $w_1$, $w_2$, \dots{}
अनंत श्रेणी नसते.

मात्र संक्रमक प्रतिलोम-सुस्थापित संबंध परिभाषित करणारे
प्रथम-क्रम सूत्र नाही. असे काही सूत्र असल्याचे
समजून, $R =
\Assign{Q}{M}$ संक्रमक आणि प्रतिलोम-सुस्थापित असेल
तेव्हा आणि केवळ तेव्हाच $\Sat{M}{\phi}$, असे माना.
नवे स्थिरांक असलेल्या विस्तारित भाषेत $\varphi_n$ हे
पुढील सूत्र असू द्या:
\[(\Atom{Q}{a_1,a_2} \land \dots \land
    \Atom{Q}{a_{n},a_{n+1}})
\]
आता पुढील सूत्रांचा संच पाहू:
\[\Gamma = \{\phi, \varphi_1, \varphi_2, \dots\}.
\]
$\Gamma$ चा प्रत्येक सांत उपसंच पूर्ततायोग्य आहे.
उपसंचात $\varphi_k$ असेल, तर $k$ हा त्यातील सर्वात मोठा
निर्देशांक घ्या; असे सूत्र नसेल, तर त्याचे मूल्य एक
घ्या. $\Domain{M_k} = \{1, \dots, k+1\}$,
आवश्यक स्थिरांकांसाठी $\Assign{a_i}{M_k} = i$,
आणि $\Assign{Q}{M_k} = <$ असे ठेवा; उर्वरित
स्थिरांकांना प्रांतातील कोणतेही मूल्य द्या.
$\{1,
\dots, k+1\}$ वरील $<$ संक्रमक आणि प्रतिलोम-सुस्थापित
असल्यामुळे $\Sat{M_k}{\phi}$. रचनेनुसार प्रत्येक
$i \le k$ साठी $\Sat{M_k}{\varphi_i}$.
प्रथम-क्रम तर्कशास्त्राच्या संहतता प्रमेयावरून
$\Gamma$ ची कोणत्यातरी संरचना~$\Struct{M}$
मध्ये पूर्ति होते. गृहीतकावरून $\Sat{M}{\phi}$
असल्यामुळे $\Assign{Q}{M}$ हा संबंध प्रतिलोम-सुस्थापित
आहे. पण $\Assign{a_1}{M}$, $\Assign{a_2}{M}$,
\dots{} यांपासून प्रतिलोम-सुस्थापितपणाने निषिद्ध
ठरवलेली अनंत श्रेणी मिळते.

दुसऱ्या दाव्याचे प्रतिउदाहरण म्हणजे वैश्विकता:
प्रत्येक $u$ आणि $v$ साठी $Ruv$. वैश्विक चौकटी
$\lforall[x][\lforall[y][\Atom{Q}{x,y}]]$ या
सूत्राने प्रथम-क्रमात परिभाषित होतात. मात्र एखादे
मोडल सूत्र नेमक्या वैश्विक चौकटींमध्येच वैध
असते असे नाही. हा निष्कर्ष स्वतंत्रपणे महत्त्वाच्या
एका तथ्यावरून येतो: वैश्विक चौकटींमध्ये वैध
असलेली सूत्रे आणि परावर्ती, सममित व संक्रमक
चौकटींमध्ये वैध असलेली सूत्रांनीमकी तीच आहेत.
परावर्ती, सममित आणि संक्रमक असूनही वैश्विक
नसलेल्या चौकटी आहेत. म्हणून सर्व वैश्विक चौकटींमध्ये
वैध असलेले प्रत्येक सूत्र काही अवैश्विक
चौकटींमध्येही वैध असते.












\section{तुल्यता संबंध आणि \Log{S5}}\label{nml:frd:es5:sec}

सर्व वैश्विक चौकटींमध्ये वैध असलेल्या सूत्रे
चा संच म्हणून मोडल तर्कशास्त्र~\Log{S5} निश्चित होते;
अशा चौकटीत प्रत्येक जग स्वतःसह प्रत्येक जगापासून
प्राप्य असते. येथे $\Box$ आवश्यकतेशी आणि
$\Diamond$ शक्यतेशी जुळतो: $\varphi$ \emph{प्रत्येक}
जगात सत्य असेल तर $\Box \varphi$ सत्य असते, आणि
$\varphi$ \emph{किमान एका} जगात सत्य असेल तर
$\Diamond \varphi$ सत्य असते. \Log{S5} हे सर्व परावर्ती,
सममित आणि संक्रमक चौकटींमध्ये, म्हणजे सर्व
\emph{तुल्यता संबंधांवर}, वैध असलेल्या सूत्रे
चा संच म्हणूनही सांगता येते.

\begin{defn}
  $W$ वरील द्विस्थानी संबंध $R$ हा \emph{तुल्यता संबंध}
  असतो तेव्हा आणि केवळ तेव्हाच, जेव्हा तो परावर्ती,
  सममित आणि संक्रमक असतो. $W$ वरील संबंध $R$
  \emph{वैश्विक} असतो तेव्हा आणि केवळ तेव्हाच, जेव्हा
  सर्व $u,v \in
  W$ साठी $Ruv$ असते.
\end{defn}

\Ax{T}, \Ax{B}, आणि \Ax{4} हे अनुक्रमे परावर्ती,
सममित आणि संक्रमक चौकटींचे वैशिष्ट्य दर्शवतात.
म्हणून प्राप्यता संबंध तुल्यता संबंध असलेल्या चौकटी
नेमक्या त्या आहेत ज्यांत ही तिन्ही सूत्रे वैध
आहेत. $R$ चे इतर परिचित गुणधर्म एकत्र केल्यास
तुल्यता संबंध परिभाषित करणाऱ्या अटींच्या समतुल्य
अटी मिळतात; म्हणून सूत्रांच्या इतर संयोगांनीही
तुल्यता संबंधांचे वैशिष्ट्य सांगता येते.

\begin{prop}\label{nml:frd:es5:prop:equivalences}
  पुढील अटी समतुल्य आहेत:
  \begin{enumerate}
  \item $R$ हा तुल्यता संबंध आहे;
  \item $R$ परावर्ती आणि युक्लिडी आहे;
  \item $R$ सर्वत्र उत्तरवर्ती, सममित आणि युक्लिडी आहे;
  \item $R$ सर्वत्र उत्तरवर्ती, सममित आणि संक्रमक आहे.
  \end{enumerate}
\end{prop}

\begin{proof}
  सरावासाठी.
\end{proof}

\begin{prob}
  पुढील विधाने दाखवून \ref{nml:frd:es5:prop:equivalences}
  सिद्ध करा:
  \begin{enumerate}
  \item $R$ सममित आणि संक्रमक असेल, तर तो युक्लिडी आहे.
  \item $R$ परावर्ती असेल, तर तो सर्वत्र उत्तरवर्ती आहे.
  \item $R$ परावर्ती आणि युक्लिडी असेल, तर तो सममित आहे.
  \item $R$ सममित आणि युक्लिडी असेल, तर तो संक्रमक आहे.
  \item $R$ सर्वत्र उत्तरवर्ती, सममित आणि संक्रमक असेल,
    तर तो परावर्ती आहे.
  \end{enumerate}
  या विधानांवरून सर्व अटी समतुल्य आहेत, हे कसे
  सिद्ध होते ते स्पष्ट करा.
\end{prob}

\ref{nml:frd:es5:prop:equivalences} हे \ref{nml:prf:prs:prop:S5}
चे चिन्हार्थविषयक प्रतिरूप आहे: $R$ तुल्यता संबंध
असलेल्या चौकटींच्या मोडल तर्कशास्त्राचे, परंपरेने
\Log{S5} म्हणतात त्याचे, ते समतुल्य वैशिष्ट्यदर्शन देते.

वैश्विक संबंध आणि तुल्यता संबंध यांचे नाते काय?
प्रत्येक वैश्विक संबंध तुल्यता संबंध असला, तरी
प्रत्येक तुल्यता संबंध वैश्विक असेलच असे नाही.
तरी सर्व वैश्विक संबंधांवर वैध असलेली सूत्रे
आणि सर्व तुल्यता संबंधांवर वैध असलेली सूत्रे
नेमकी तीच आहेत.

\begin{prop}
  $R$ हा तुल्यता संबंध असू द्या. प्रत्येक $w \in W$
  साठी $w$ चा \emph{तुल्यतावर्ग} म्हणजे
  $[w] = \{w'\in W :
  Rww'\}$ हा संच ठरवा. मग:
  \begin{enumerate}
  \item $w \in [w]$;
  \item प्रत्येक तुल्यतावर्ग~$[w]$ वर $R$ वैश्विक आहे;
  \item तुल्यतावर्ग $W$ ला परस्पर विसंयुक्त आणि
    एकत्रितपणे संपूर्ण W व्यापणाऱ्या उपसंचांत विभागतात.
  \end{enumerate}
\end{prop}

\begin{prop}\label{nml:frd:es5:prop:S5=univ}
  सूत्र~$\varphi$ हे $R$ तुल्यता संबंध असलेल्या
  सर्व चौकटींमध्ये $\mModel{F} =\tuple{W,R}$ वैध
  असते तेव्हा आणि केवळ तेव्हाच, जेव्हा ते $R$
  वैश्विक असलेल्या सर्व चौकटींमध्ये
  $\mModel{F} =\tuple{W,R}$ वैध असते. म्हणून
  वैश्विक चौकटींचे तर्कशास्त्र म्हणजे \Log{S5}.
\end{prop}

\begin{proof}
  $W$ वरील वैश्विक संबंध~$R$ हा तुल्यता संबंध असतो,
  हे लगेच तपासता येते. म्हणून $\varphi$ हे $R$ तुल्यता
  संबंध असलेल्या सर्व चौकटींमध्ये वैध असेल, तर ते
  सर्व वैश्विक चौकटींमध्येही वैध असते. उलट दिशेसाठी
  प्रतिपरिवर्तनाने युक्तिवाद करू. $W$ वरील $R$
  तुल्यता संबंध असलेल्या चौकट~$\tuple{W,R}$ वर
  आधारित $\mModel{M} = \tuple{W, R, V}$ या
  प्रतिमानात $\psi$ हे सूत्र एखाद्या जगात~$w$
  असत्य आहे असे समजा. म्हणजे $\mSat/{M}{\psi}[w]$.
  पुढीलप्रमाणे $\mModel{M}' = \tuple{W',
    R', V'}$ हे प्रतिमान ठरवा:
  \begin{enumerate}
  \item $W' = [w]$;
  \item $R'$ हा $W'$ वर वैश्विक आहे;
  \item $V'(p) = V(p) \cap W'$.
  \end{enumerate}
  (\ref{nml:frd:es5:fig:partition} मधील छायांकित भाग
  $\mModel{M}'$ चा जगांचा संच~$W'$ दाखवतो.)
  $W'$ वर $R$ आणि $R'$ हे संबंध जुळतात, हे सहज
  दिसते. मग सूत्रे वर विगमन करून, प्रत्येक
  $\varphi$ आणि सर्व $w' \in W'$ साठी
  $\mSat{M'}{\varphi}[w']$ तेव्हा आणि केवळ तेव्हाच,
  जेव्हा $\mSat{M}{\varphi}[w']$, हे दाखवता येते
  (कारण $W'
  \subseteq W$). विशेषतः $\mSat/{M'}{\psi}[w]$;
  म्हणून $\psi$ हे वैश्विक चौकटीवर आधारित
  प्रतिमानात असत्य आहे.
\end{proof}

\begin{figure}[t]
  \centering
  \begin{tikzpicture}[node distance=2cm, auto, thick]
    \clip [rounded corners] (0,0) -- (6,0) -- (6,4) -- (0,4) --  cycle;
    \begin{scope}
      \clip (0,0) -- (2,0)  .. controls (1.5,1.5) and (4.5,1.5)
      .. (4,4) -- (0,4) -- (0,0) -- cycle;
      \filldraw[fill=gray!40] (0,4) -- (0,2) .. controls (1.5,1.5)
      and (4.5,1.5) .. (4.5,0) -- (6,0) -- (6,4) -- (0,4) --  cycle;
    \end{scope}
    \draw [name path=line1] (2,0) .. controls (1.5,1.5) and (4.5,1.5) .. (4,4);
    \draw [name path=line2] (0,2)  .. controls (1.5,1.5) and (4.5,1.5) .. (4.5,0);
    \path [name intersections={of=line1 and line2}];
    \draw [rounded corners, use as bounding box, very thick] (0,0)
    -- (6,0) -- (6,4) -- (0,4) --  cycle;
    \node at (2,3) {$[w]$} ;
    \node at (1,1) {$[u]$} ;
    \node at (3,0.75) {$[v]$} ;
    \node at (4.75,2) {$[z]$} ;
  \end{tikzpicture}
  \caption{$W$ चा तुल्यतावर्गांमध्ये विभाग.}
  \label{nml:frd:es5:fig:partition}
\end{figure}













\section{द्वितीय-क्रमातील परिभाष्यता}\label{nml:frd:st:sec}

मोडल सूत्रांनी परिभाषित होणारा चौकटींचा प्रत्येक गुणधर्म
प्रथम-क्रमात परिभाष्य असेलच असे नाही. पण एकस्थानी
विधेयांवर संख्यापन करण्याची परवानगी दिल्यास (म्हणजे
एकस्थानी द्वितीय-क्रम संख्यापन), मोडल सूत्रांनी परिभाष्य
असलेले सर्व चौकट-गुणधर्म परिभाषित करता येतात.
एखादे मोडल सूत्र जगात सत्य असण्याच्या अटी
प्रथम-क्रम सूत्रे शी ज्या पद्धतशीर रीतीने
संबंधित असतात, तिचा उपयोग करणे ही युक्ती आहे.
या पद्धतीला मोडल सूत्रांचे प्रथम-क्रम सूत्रांत
मानक भाषांतर म्हणतात. त्यासाठीच्या भाषेत
प्राप्यता संबंधासाठी द्विस्थानी विधेय~$Q$
तर असतेच, शिवाय~$\varphi$ मध्ये येणाऱ्या
विधानीय चले~$\Obj p_i$ साठी
एकस्थानी विधेय~$P_i$ देखील असते.

\begin{defn}
  \emph{मानक भाषांतर}~$\ST_x(\varphi)$ ची विगमनाने
  पुढीलप्रमाणे व्याख्या करतो. मोडल उपवाक्यांत $x$ पेक्षा वेगळा
  नवीन व्यक्तिचल $y$ निवडा. आतील भाषांतरातील संख्यापकांसाठीही
  मुक्त चल बद्ध होणार नाहीत अशी नवीन बद्ध चलांची नावे निवडावीत:
  \begin{enumerate}
  \item \indcase{\varphi}{\lfalse}{$\ST_x(\indfrm) = \lfalse$.}
  \item \indcase{\varphi}{\ltrue}{$\ST_x(\indfrm) = \ltrue$.}
  \item \indcase{\varphi}{\Obj p_i}{$\ST_x(\indfrm) = \Atom{P_i}{x}$.}
  \item \indcase{\varphi}{\lnot \psi}{$\ST_x(\indfrm) = \lnot \ST_x(\psi)$.}
  \item \indcase{\varphi}{(\psi \land \chi)}{$\ST_x(\indfrm) = (\ST_x(\psi)
      \land \ST_x(\chi))$.}
  \item \indcase{\varphi}{(\psi \lor \chi)}{$\ST_x(\indfrm) = (\ST_x(\psi)
      \lor \ST_x(\chi))$.}
  \item \indcase{\varphi}{(\psi \lif \chi)}{$\ST_x(\indfrm) = (\ST_x(\psi)
      \lif \ST_x(\chi))$.}
  \item \indcase{\varphi}{(\psi \liff \chi)}{$\ST_x(\indfrm) = (\ST_x(\psi)
      \liff \ST_x(\chi))$.}
  \item \indcase{\varphi}{\Box \psi}{$\ST_x(\indfrm) =
      \lforall[y][(\Atom{Q}{x,y} \lif \ST_y(\psi))]$.}
  \item \indcase{\varphi}{\Diamond \psi}{$\ST_x(\indfrm) =
      \lexists[y][(\Atom{Q}{x,y} \land \ST_y(\psi))]$.}
  \end{enumerate}
\end{defn}

उदाहरणार्थ, $\ST_x(\Box p \lif p)$ हे
$\lforall[y][(\Atom{Q}{x,y}
  \lif \Atom{P}{y})] \lif \Atom{P}{x}$ असे आहे.
~$\ST_x(\varphi)$ च्या भाषेतील कोणत्याही संरचना
साठी प्रांत, $Q$ ला नेमून दिलेला द्विस्थानी संबंध,
आणि एकस्थानी विधेये~$P_i$ यांना नेमून
दिलेले प्रांताचे उपसंच लागतात. दुसऱ्या शब्दांत,
अशा संरचनेचे घटक हे~$\varphi$ च्या प्रतिमानाच्या
घटकांशी तंतोतंत जुळतात: प्रांत म्हणजे जगांचा संच;
$Q$ ला नेमलेला द्विस्थानी संबंध म्हणजे प्राप्यता
संबंध; आणि $P_i$ यांना नेमलेले उपसंच म्हणजे
नेमणुका~$V(\Obj p_i)$. त्यामुळे मोडल प्रतिमानातील
$\varphi$ ची पूर्ति आणि अनुरूप संरचना मधील
$\ST_x(\varphi)$ ची पूर्ति जुळते, हे अपेक्षितच आहे:

\begin{prop}\label{nml:frd:st:prop:st}
  $\mModel{M} = \tuple{W, R, V}$ असू द्या.
  $\Domain{M'} = W$, $\Assign{Q}{M'} = R$, आणि
  $\Assign{P_i}{M'} = V(\Obj p_i)$ असलेली
  $\Struct{M'}$ ही प्रथम-क्रम संरचना असू द्या,
  तसेच $s(x) = w$ असू द्या. मग
  \[  \mSat{M}{\varphi}[w] \text{ तेव्हा आणि केवळ तेव्हाच } \Sat{M'}{\ST_x(\varphi)}[s]
  \]
\end{prop}

\begin{proof}
  $\varphi$ वर विगमनाने.
\end{proof}

\begin{prop}
  $\varphi$ हे मोडल सूत्र आणि
  $\mModel{F} = \tuple{W, R}$ ही चौकट असू द्या.
  $\Domain{F'} = W$ आणि $\Assign{Q}{F'} = R$
  असलेली $\Struct{F'}$ ही प्रथम-क्रम संरचना
  असू द्या. पुढील द्वितीय-क्रम सूत्राला $\varphi'$
  म्हणा:
  \[  \lforall[X_1][\dots\lforall[X_n][\lforall[x][
        \SSubst{\ST_x(\varphi)}{\subst{X_1}{P_1}, \dots,
          \subst{X_n}{P_n}}]]],
  \]
  जिथे $P_1$, \dots, $P_n$ ही~$\ST_x(\varphi)$
  मधील सर्व एकस्थानी विधेये आहेत. मग
  \[  \mModel{F} \Entails \varphi \text{ तेव्हा आणि केवळ तेव्हाच } \Sat{F'}{\varphi'}
  \]
\end{prop}

\begin{proof}
  $\Sat{F'}{\varphi'}$ तेव्हा आणि केवळ तेव्हाच, जेव्हा
  $i = 1$, \dots,~$n$ साठी
  $\Assign{P_i}{M'} \subseteq W$ असलेल्या
  मूळ रचनेच्या प्रत्येक विस्ताररूप
  संरचना~$\mModel{M'}$ मध्ये, आणि
  $s(x) \in W$ असलेल्या प्रत्येक $s$ साठी,
  $\Sat{M'}{\ST_x(\varphi)}[s]$ असते.
  \ref{nml:frd:st:prop:st} नुसार, हे तेव्हा आणि केवळ
  तेव्हाच, जेव्हा $\mModel{F}$ वर आधारित सर्व
  प्रतिमाने~$\mModel{M}$ आणि प्रत्येक जग~$w \in W$
  यांसाठी $\mSat{M}{\varphi}[w]$ असते; म्हणजे
  $\mModel{F} \Entails \varphi$.
\end{proof}

\begin{defn}
  चौकटींचा वर्ग~$\mClass{F}$ \emph{द्वितीय-क्रमात
  परिभाष्य} असतो, जर एकच द्विस्थानी
  विधेय~$P$ असलेल्या आणि द्वितीय-क्रमाचे संख्यापक
  केवळ एकस्थानी संचचलांवर असलेल्या भाषेत असे
  वाक्य~$\varphi$ असेल की,
  $\Domain{M} = W$ आणि $\Assign{P}{M} = R$
  असलेल्या संरचना~$\Struct{M}$ मध्ये
  $\Sat{M}{\varphi}$ तेव्हा आणि केवळ तेव्हाच,
  जेव्हा $\mModel{F} = \tuple{W, R} \in \mClass{F}$.
  व्यक्तिचलांवरील प्रथम-क्रम संख्यापकही अनुमत आहेत.
\end{defn}

\begin{cor}
  चौकटींचा वर्ग एखाद्या सूत्र~$\varphi$ ने
  परिभाष्य असेल, तर त्याच्या अनुरूप प्राप्यता
  संबंधांचा वर्ग एकस्थानी द्वितीय-क्रम वाक्याने
  परिभाष्य असतो.
\end{cor}

\begin{proof}
  आधीच्या सिद्धतेतील एकस्थानी द्वितीय-क्रम
  वाक्य~$\varphi'$ ला आवश्यक गुणधर्म आहे.
\end{proof}

उदाहरण म्हणून पुन्हा सूत्र~$\Box p \lif p$
पाहा. ते परावर्तितेची व्याख्या करते. अर्थात,
परावर्तिता वाक्य~$\lforall[x][\Atom{Q}{x,x}]$
याने प्रथम-क्रमात परिभाष्य आहे. पण ती पुढील
एकस्थानी द्वितीय-क्रम वाक्य नेही
परिभाष्य आहे:
\[\lforall[X][\lforall[x][(\lforall[y][(\Atom{Q}{x,y} \lif \Atom{X}{y})]
    \lif \Atom{X}{x})]].
\]
म्हणजे ही दोन वाक्ये समतुल्य आहेत. आधी द्वितीय-क्रम वाक्य
संरचना~$\Struct{M}$ मध्ये सत्य आहे असे समजा.
कोणतेही $w \in W$ घ्या. $R=\Assign{Q}{M}$ असू द्या आणि
$S=\Setabs{z \in W}{Rwz}$ ठेवा. $s(x)=w$ आणि $s(X)=S$
असलेली नेमणूक $s$ निवडा. वैश्विक संख्यापनाच्या व्याख्येवरून
\[\Sat{M}{\lforall[y][(\Atom{Q}{x,y} \lif \Atom{X}{y})]
    \lif \Atom{X}{x}}[s].
\]
$S$ च्या निवडीवरून या अभिव्यंजनाचे पूर्वांग $s$ नुसार सत्य आहे:
$w$ पासून प्राप्य असलेले प्रत्येक जग $S$ मध्ये आहे.
म्हणून $\Sat{M}{\Atom{X}{x}}[s]$; म्हणजे $w \in S$.
$S$ च्या व्याख्येनुसार $Rww$. $w$ कोणतेही असल्यामुळे
$\Sat{M}{\lforall[x][\Atom{Q}{x,x}]}$ मिळते.

आता $\Sat{M}{\lforall[x][\Atom{Q}{x,x}]}$
असे समजा आणि
$\Sat{M}{\lforall[X][\lforall[x][(\lforall[y][(\Atom{Q}{x,y} \lif
        \Atom{X}{y})] \lif \Atom{X}{x})]]}$
हे दाखवू. कोणतीही नेमणूक~$s$ घ्या आणि
$\Sat{M}{\lforall[y][(\Atom{Q}{x,y} \lif
    \Atom{X}{y})]}[s]$ असे माना. $s'$ ही~$s$
ची $y$-भिन्नरूप नेमणूक असू द्या, जिथे $s'(y)
= s(x)$; मग
$\Sat{M}{\Atom{Q}{x,y} \lif \Atom{X}{y}}[s']$,
म्हणजे $\Sat{M}{\Atom{Q}{x,x} \lif \Atom{X}{x}}[s]$.
$\Sat{M}{\lforall[x][\Atom{Q}{x,x}]}$ असल्यामुळे
पूर्वांग सत्य आहे; म्हणून $\Sat{M}{\Atom{X}{x}}[s]$,
आणि हेच दाखवायचे होते.

चौकटींचे काही परिभाष्य वर्ग प्रथम-क्रमात परिभाष्य
नसल्यामुळे, $\varphi'$ या रूपातील प्रत्येक एकस्थानी
द्वितीय-क्रम वाक्य हे प्रथम-क्रम
वाक्य शी समतुल्य असेलच असे नाही.
कोणती अशी आहेत हे ठरवण्याची प्रभावी पद्धत नाही.



